\section{可定向性}\label{sec:orientability}

设 $V$ 是有限维向量空间, 考虑两组有序基 $\beta=\{b_1,\ldots,b_n\}$ 和 $\beta'=\{b'_1,\ldots,b'_n\}$. 存在唯一的线性变换 $S:V\to V$, 使得对每个 $i=1,\ldots,n$ 都有 $b'_i=Sb_i$. 若 $\det S>0$, 就称这两组基\textbf{等价}. 这样定义了一个等价关系, 把 $V$ 的所有有序基分成两个等价类. $V$ 的一个\textbf{定向}是给一个等价类的元素赋予正号, 给另一个等价类的元素赋予负号. 赋给一组基的符号称为它的\textbf{定向}; 根据其符号, 称这组基为\textbf{正定向}或\textbf{负定向}. 显然, $V$ 恰有两种可能的定向.

\begin{RiemannRemark}\label{rem:orders-and-standard-orientation}
\begin{enumerate}[label=(\arabic*),leftmargin=*]
\item 基的次序很重要. 交换两个基向量的位置, 就得到另一组有序基, 其定向与原来相反.
\item 零维向量空间的一个定向就是赋予符号 $+1$ 或 $-1$.
\item 将正号赋给标准有序基的定向称为 $\mathbb{R}^n$ 的\textbf{标准定向}.
\end{enumerate}
\end{RiemannRemark}

两个已定向向量空间之间的同构 $A:V\to W$ 将 $V$ 的等价有序基映为 $W$ 的等价有序基. 因此, 对任何有序基 $\beta$, 像 $A\beta$ 的符号要么总与 $\beta$ 的符号相同, 要么总与之相反. 第一种情形下, 称同构 $A$ \textbf{保定向}; 第二种情形下, 称其\textbf{反定向}.

光滑流形的一个\textbf{定向}是对所有切空间 $T_pM$ 选取定向. 若 $\dim M=n\geq 1$, 这些定向必须光滑地相容, 即对每个点 $p\in M$, 存在 $p$ 附近的参数化 $(U,\varphi)$, 使得
\[
  (\mathrm{d}\varphi)_x:\mathbb{R}^n\to T_{\varphi(x)}M
\]
在每个点 $x\in U$ 处保持 $\mathbb{R}^n$ 的标准定向.

\begin{RiemannRemark}\label{rem:orientation-in-dimension-zero}
若 $M$ 的维数为零, 则一个定向就是给每个点 $p\in M$ 赋予符号 ($+1$ 或 $-1$), 称为\textbf{定向数}.
\end{RiemannRemark}

\begin{RiemannDefinition}\label{def:orientable-manifold}
若光滑流形 $M$ 容许一个定向, 则称 $M$ \textbf{可定向}.
\end{RiemannDefinition}

\begin{RiemannProposition}\label{prop:two-orientations-on-connected-manifold}
若光滑流形 $M$ 连通且可定向, 则它恰有两个定向.
\end{RiemannProposition}

\begin{RiemannProof}
我们证明, 两个定向一致的点所成的集合与两者不一致的点所成的集合都是开集. 因此, 其中一个必为 $M$, 另一个必为空集. 设 $p$ 是 $M$ 中的一点, $(U_\alpha,\varphi_\alpha)$ 和 $(U_\beta,\varphi_\beta)$ 是以 $p$ 为中心的两个参数化, 使得 $\mathrm{d}\varphi_\alpha$ 对第一个定向保定向, 而 $\mathrm{d}\varphi_\beta$ 对第二个定向保定向. 映射 $\bigl(\mathrm{d}(\varphi_\beta^{-1}\circ\varphi_\alpha)\bigr)_0:\mathbb{R}^n\to\mathbb{R}^n$ 或者保定向 (如果两个定向在 $p$ 处一致), 或者反定向. 在第一种情形下, 它在 $0$ 处的行列式为正, 因而由连续性, 对 $0$ 的某个邻域内的 $x$, $\bigl(\mathrm{d}(\varphi_\beta^{-1}\circ\varphi_\alpha)\bigr)_x$ 的行列式为正, 这意味着两个定向在 $p$ 的一个邻域内一致. 类似地, 若 $\bigl(\mathrm{d}(\varphi_\beta^{-1}\circ\varphi_\alpha)\bigr)_0$ 反定向, 则 $\bigl(\mathrm{d}(\varphi_\beta^{-1}\circ\varphi_\alpha)\bigr)_x$ 的行列式在 $0$ 的一个邻域内为负, 所以两个定向在 $p$ 的一个邻域内不一致.

设 $O$ 是 $M$ 的一个定向 (即对每个 $p\in M$, 光滑地选取 $T_pM$ 的定向 $O_p$), $-O$ 是相反的定向 (对应于在每个切空间 $T_pM$ 上取相反的定向 $-O_p$). 若 $O'$ 是 $M$ 的另一个定向, 则对给定的点 $p\in M$, $O'_p$ 要么与 $O_p$ 一致, 要么与 $-O_p$ 一致 (因为向量空间只有两种可能的定向). 因此, $O'$ 在 $M$ 上或者与 $O$ 一致, 或者与 $-O$ 一致.
\end{RiemannProof}

下述命题给出可定向性的另一种刻画.

\begin{RiemannProposition}\label{prop:orientation-preserving-atlas}
光滑流形 $M$ 可定向, 当且仅当存在图册 $\mathcal{A}=\{(U_\alpha,\varphi_\alpha)\}$, 使所有重叠映射 $\varphi_\beta^{-1}\circ\varphi_\alpha$ 都保定向.
\end{RiemannProposition}

\begin{RiemannProof}
Exercise~\ref{ex:orientability}(\ref{ex:characterize-orientable-manifold-by-atlas}).
\end{RiemannProof}

\textbf{已定向流形}是带有一个选定定向的可定向流形. 对两个同维的已定向流形之间的映射 $f:M\to N$, 若 $(\mathrm{d}f)_p$ 在所有点 $p\in M$ 处都保定向, 则称 $f$ \textbf{保定向}; 若 $(\mathrm{d}f)_p$ 在所有点 $p\in M$ 处都反定向, 则称 $f$ \textbf{反定向}.

\begin{RiemannExercise}\label{ex:orientability}
\begin{enumerate}[label=(\arabic*),ref=\arabic*,leftmargin=*]
\item\label{ex:equivalence-of-ordered-bases} 证明上述有限维向量空间的有序基之间的"等价"关系确实是等价关系.
\item\label{ex:characterize-orientable-manifold-by-atlas} 证明可微流形 $M$ 可定向, 当且仅当存在图册 $\mathcal{A}=\{(U_\alpha,\varphi_\alpha)\}$, 使所有重叠映射 $\varphi_\beta^{-1}\circ\varphi_\alpha$ 都保定向.
\item\label{ex:orientability-of-sphere-and-two-chart-manifolds}
\begin{enumerate}[label=(\alph*),leftmargin=*]
\item 证明若流形 $M$ 可由两个坐标邻域 $V_1,V_2$ 覆盖, 且 $V_1\cap V_2$ 连通, 则 $M$ 可定向.
\item 证明 $S^n$ 可定向.
\end{enumerate}
\item\label{ex:orientation-of-bases-along-curve} 设 $M$ 是已定向的 $n$ 维流形, $c:I\to M$ 是可微曲线. \textbf{沿 $c$ 的光滑向量场}是可微映射 $V:I\to TM$, 使得对所有 $t\in I$ 有 $V(t)\in T_{c(t)}M$ (参见第~\RGref{chap:riemannian-manifolds}~章第~\RGref{sec:affine-connections}~节). 证明若 $V_1,\ldots,V_n:I\to M$ 是沿 $c$ 的光滑向量场, 且对所有 $t\in I$, $\{V_1(t),\ldots,V_n(t)\}$ 都是 $T_{c(t)}M$ 的一组基, 则这些基都有相同的定向.
\item\label{ex:mobius-band-is-not-orientable} 可以将 Möbius 带视为 $\mathbb{R}^3$ 中由浸入 $g:(-1,1)\times\mathbb{R}\to\mathbb{R}^3$ 的像给出的 2 维子流形, 其中
\[
  g(t,\varphi)=\left(
    \left(1+t\cos\left(\frac{\varphi}{2}\right)\right)\cos\varphi,
    \left(1+t\cos\left(\frac{\varphi}{2}\right)\right)\sin\varphi,
    t\sin\left(\frac{\varphi}{2}\right)\right).
\]
证明 Möbius 带不可定向.
\item\label{ex:orientability-preserved-by-diffeomorphism} 设 $f:M\to N$ 是两个光滑流形之间的微分同胚. 证明 $M$ 可定向当且仅当 $N$ 可定向. 若另外两个流形都连通且已定向, 并且在某一点 $p\in M$ 处 $(\mathrm{d}f)_p:T_pM\to T_{f(p)}N$ 保定向, 证明 $f$ 保定向.
\item\label{ex:well-defined-product-orientation} 设 $M$ 和 $N$ 是两个已定向流形. 按如下方式定义乘积流形 $M\times N$ 的定向 (称为\textbf{乘积定向}): 若 $\alpha=\{a_1,\ldots,a_m\}$ 和 $\beta=\{b_1,\ldots,b_n\}$ 分别为 $T_pM$ 和 $T_qN$ 的有序基, 考虑 $T_{(p,q)}(M\times N)\cong T_pM\times T_qN$ 的有序基 $\{(a_1,0),\ldots,(a_m,0),(0,b_1),\ldots,(0,b_n)\}$. 令这组基的符号等于 $\alpha$ 与 $\beta$ 的符号之积, 以此定义该空间的一个定向. 证明这个定向与 $\alpha$ 和 $\beta$ 的选取无关.
\item\label{ex:orientability-of-tangent-bundle} 证明切丛 $TM$ 总是可定向的, 即使 $M$ 不可定向.
\item\label{ex:orientable-double-covering} (\emph{可定向二重覆盖}) 设 $M$ 是不可定向的 $n$ 维流形. 对每个点 $p\in M$, 考虑 $T_pM$ 的基的 (两个) 等价类所成的集合 $\mathcal{O}_p$. 设 $\overline M$ 为集合
\[
  \overline M=\{(p,O_p)\mid p\in M,\ O_p\in\mathcal{O}_p\}.
\]
给定 $M$ 的一个参数化 $(U,\varphi)$, 考虑由下式定义的映射 $\overline{\varphi}:U\to\overline M$:
\[
  \overline{\varphi}(x^1,\ldots,x^n)
  =\left(\varphi(x^1,\ldots,x^n),
    \left[\left(\frac{\partial}{\partial x^1}\right)_{\varphi(x)},
       \ldots,\left(\frac{\partial}{\partial x^n}\right)_{\varphi(x)}\right]\right),
\]
其中 $x=(x^1,\ldots,x^n)\in U$, 且 $\bigl[\bigl(\frac{\partial}{\partial x^1}\bigr)_{\varphi(x)},\allowbreak\ldots,\bigl(\frac{\partial}{\partial x^n}\bigr)_{\varphi(x)}\bigr]$ 表示 $T_{\varphi(x)}M$ 的基 $\bigl\{\bigl(\frac{\partial}{\partial x^1}\bigr)_{\varphi(x)},\allowbreak\ldots,\bigl(\frac{\partial}{\partial x^n}\bigr)_{\varphi(x)}\bigr\}$ 的等价类.
\begin{enumerate}[label=(\alph*),leftmargin=*]
\item 证明这些映射在 $\overline M$ 上确定了一个 $n$ 维可定向可微流形的结构.
\item 考虑由 $\pi(p,O_p)=p$ 定义的映射 $\pi:\overline M\to M$. 证明 $\pi$ 可微且满射. 此外, 证明对每个 $p\in M$, 存在 $p$ 的一个邻域 $V$, 使得 $\pi^{-1}(V)=W_1\cup W_2$, 其中 $W_1$ 和 $W_2$ 是 $\overline M$ 的两个不相交的开子集, 且 $\pi$ 到 $W_i$ ($i=1,2$) 的限制是到 $V$ 的微分同胚.
\item 证明 $\overline M$ 连通 (因此称 $\overline M$ 是 $M$ 的\textbf{可定向二重覆盖}).
\item 设 $\sigma:\overline M\to\overline M$ 由 $\sigma(p,O_p)=(p,-O_p)$ 定义, 其中 $-O_p$ 表示 $T_pM$ 与 $O_p$ 相反的定向. 证明 $\sigma$ 是反定向的微分同胚, 且满足 $\pi\circ\sigma=\pi$ 和 $\sigma\circ\sigma=\mathrm{id}$.
\item 证明任何单连通流形都是可定向的.
\end{enumerate}
\end{enumerate}
\end{RiemannExercise}
